\section{Conclusion}\label{sec:conclusion}

We measured how much of a retrieval gain reaches the answer when a RAG system swaps one open embedding model for another. On single-hop NQ, between \ev{T.lnq.R4} and \ev{T.lnq.R3} EM points followed each Hit@5 point that a dense encoder gained over MiniLM-L6, for both readers (RQ1); against BM25 the ratios were \ev{V4.b0.nq.min} to \ev{V4.b0.nq.max}. On multi-hop HotpotQA the Hit@5-based ratio was larger; in a post hoc decomposition, \ev{DS.hp.min} to \ev{DS.hp.max} of it came from higher EM given a hit, in line with better encoders retrieving the second supporting passage more often, which Hit@5 does not count; with Both@5 as the denominator, the HotpotQA ratios were \ev{V4.both.hp.min} to \ev{V4.both.hp.max}. Among dense encoders, bge-large gave a small EM gain over e5-base-v2 on HotpotQA and none detectable on NQ, and Qwen3-Embedding-0.6B gave lower EM on HotpotQA (RQ2). Top-ranked non-gold passages lowered EM, gold position did not change it detectably, a reranker helped only the weak BM25 retrieval stage, and ten passages narrowed the gaps between retrievers (RQ3). Encoding throughput differed \ev{E7.ratio.R1R6}-fold, while approximate search returned nearly the same top passages and the four retrievers tested on the full NQ corpus kept their order (RQ4).

Within this setup, retrieval gains overstated answer gains on single-hop questions, and on multi-hop questions once retrieval was scored by Both@5; Hit@5 hid much of the multi-hop retrieval gain. We recommend that RAG studies report answer gains next to retrieval gains, with paired intervals and the baseline stated, and use retrieval metrics that require all evidence a question needs. Future work should test larger readers and other domains, run the readers on the full corpus and on approximate-search results, and replace EM with validated human or model judgements.
